\documentclass[conference]{IEEEtran}
\IEEEoverridecommandlockouts

\usepackage{cite}
\usepackage{amsmath,amssymb,amsfonts}
\usepackage{graphicx}
\usepackage{textcomp}
\usepackage{xcolor}
\usepackage{booktabs}
\usepackage{microtype}
\usepackage[colorlinks=true,allcolors=blue!55!black]{hyperref}

\newcommand{\code}[1]{\texttt{\small #1}}
\newcommand{\cls}[1]{\textsc{#1}}

\begin{document}

\title{Finding BOOM CPU Bugs with an AI Loop\\
and Two Reference Models}

\author{\IEEEauthorblockN{Divya Kohli}
\IEEEauthorblockA{\textit{University of California, Santa Cruz}\\
dkohli1@ucsc.edu}
\thanks{A3 CHIA Hackathon 2026, MICRO A3 Workshop.}}

\maketitle

\begin{abstract}
We built a tool that looks for bugs in the BOOM CPU. BOOM is an out-of-order
RISC-V core. The tool runs a test program on BOOM and also on reference models
that follow the RISC-V rules, then compares the results. If BOOM does something
different, that may be a bug. The tool has four parts. The first part reads
BOOM's source code and picks spots that often have bugs. The second part writes
small test programs aimed at those spots. The third part runs each program on
BOOM and on the reference models and compares what they commit. The fourth part
looks at each difference and decides what it means: a real BOOM bug, a bug in a
reference model, a difference the rules allow, or a problem with the test setup.
We use two reference models, Spike and Dromajo, not one. With two, the tool can
often decide by majority vote instead of guessing. To measure how well it works,
we planted known bugs in BOOM's source and checked how many the tool finds. It
found every planted bug that changes what the CPU commits, and it never reported
a bug on the clean version. We are clear about what we built and what already
existed in the Berkeley tools.
\end{abstract}

\begin{IEEEkeywords}
CPU verification, RISC-V, BOOM, differential testing, co-simulation, CHIA,
hardware/software co-design
\end{IEEEkeywords}

\section{Introduction}
Comparing a real CPU design against a simple reference model is a common way to
find CPU bugs. The Chipyard tools~\cite{chipyard} already let you run Spike next
to a BOOM simulation and compare them. So running the comparison is not the hard
part. The hard part is deciding what to test, writing tests that reach the
tricky cases, and, most of all, deciding what a difference means. When BOOM and
a reference model disagree, it can be a bug in BOOM, a bug in the reference
model, a difference the rules allow, or just a problem with how we ran the test.

This work builds a loop around that last problem. We are honest about credit:
running Spike next to BOOM is an existing Berkeley feature, and we use it as is.
What we add is the layer on top: reading the RTL to choose where to test,
writing tests aimed there, and deciding what each difference means with a strong
rule against calling something a BOOM bug when we are not sure.

\section{Background}
\textbf{CHIA}~\cite{chia} is an open framework for AI-driven hardware and
software co-design flows. It already supports many tools, including
Chipyard, Spike~\cite{spike}, Verilator~\cite{verilator}, and gem5~\cite{gem5}.
\textbf{BOOM}~\cite{boom} is an out-of-order RISC-V~\cite{riscv} core; we test
its third generation, SonicBOOM~\cite{sonicboom}, which is written in the Chisel
hardware language~\cite{chisel}. A
\textbf{reference model} runs instructions one at a time and computes the
correct result directly. An out-of-order core does the same work in a more
complex way, with renaming, guessing ahead, queues, and forwarding. Those are
the places where it can go wrong. Because a reference model has no idea of
clock cycles, we only compare \emph{committed state}: the list of instructions
that retired, their bits, and the register values they wrote. Cycle counts and
issue order are not compared.

\section{Related Work}
The closest work is CPU fuzzing. RFUZZ~\cite{rfuzz} fuzzes RTL on FPGAs.
DIFUZZRTL~\cite{difuzzrtl} adds register-coverage guidance and differential
testing. TheHuzz~\cite{thehuzz} fuzzes at the instruction level against a golden
model. ProcessorFuzz~\cite{processorfuzz} guides the search using changes in
control and status registers and tests Rocket, BOOM, and BlackParrot.
Cascade~\cite{cascade} writes long, tangled programs that stop running when they
hit a bug. Constrained-random generators like riscv-dv~\cite{riscvdv} also make
legal test programs. All of these focus on \emph{making test inputs and
measuring coverage}. Our loop is different in two ways. First, it chooses where
to test by reading the RTL and pointing tests at named spots, not by mutating
random inputs. Second, and more important, we treat \emph{deciding what a
difference means} as the main problem, not finding the difference. This matters:
ProcessorFuzz found a bug in a reference model itself, which shows a difference
is not always a bug in the design. The classic idea of comparing two
implementations of the same rules comes from compiler testing, such as
Csmith~\cite{csmith}. Separately, large language models are now used to write
and check hardware: ChipNeMo~\cite{chipnemo} adapts a model to chip-design tasks,
and VerilogEval~\cite{verilogeval} benchmarks LLMs on Verilog generation. Our
work uses a model to choose and write tests, not to write the design.

\section{How the Loop Works}
The loop has four parts. Each is a separate CHIA block.

\subsection{Pick spots}
This part reads BOOM's source and looks for spots that often hold bugs, in five
groups: exception priority, queue pointers, forwarding and bypass, memory
ordering, and code the authors marked as unfinished (\code{TODO} and disabled
features). Each hit is saved with a \code{file:line} pointer, so every later
test can be traced back to the exact line that motivated it. On BOOM v3 (56
Scala files, about 18.8k lines) the scan finds 423 spots. A model then turns the
top spots into test ideas, each with a setup, a test sketch, and the exact
difference it expects to see. We rank the spots so a memory-ordering idea points
at the load/store unit, not at an unrelated line that only matched the same
keyword.

\subsection{Write tests}
This part turns each idea into a small bare-metal RISC-V program. It wraps the
program with a trap handler and a clean exit. The key part is a check: a test is
only kept if it assembles, finishes, and runs enough instructions to do real
work. A program that runs nothing produces a clean comparison and proves
nothing, so a loop that counts clean runs as success would happily report that a
broken test writer found no bugs. Rejected tests are counted and reported. They
mean the check is working, not that the loop failed.

\subsection{Run and compare}
This part turns both models' output into one common list of committed
instructions and compares them step by step from a start point, stopping at the
first difference. It starts from the program's entry, not from reset, because
the two models boot differently. A useful practical note: BOOM's commit log uses
almost the same format as Spike, so one parser reads both.

\subsection{Decide what it means}
This part sorts each difference into one of five kinds: a BOOM bug, a
reference-model bug, an allowed difference, a test-setup problem, or not sure.
Easy cases are decided first with no model call. A difference that does not
repeat across runs of two fixed models is not sure. A difference on a
free-running counter (\code{cycle}, \code{time}, \code{instret}) is allowed by
the rules. An empty or cut-off trace is a test-setup problem, and we tell the
two apart because saying ``raise the instruction budget'' is useless when the
simulator never started. Only what is left goes to the model, and with real
evidence: the disassembly at both program counters, nearby trace, and the RTL
line the idea pointed at. The bar for calling something a BOOM bug is high on
purpose. A wrong bug report to BOOM's authors costs more than a missed one, so
anything the evidence does not settle is marked ``not sure'' instead of being
rounded up. Only real BOOM-bug and reference-bug results feed back into the next
round of ideas.

\subsection{Vote with two reference models}
Older CPU fuzzers use one reference model, so a difference leaves ``which side
is wrong?'' to a person. We add a second reference model, Dromajo~\cite{dromajo},
next to Spike, and let the tool vote. Call the implementation $B$ and the two
references $A_1$ and $A_2$. If all three agree, it is clean. If both references
agree and $B$ is the odd one out, it is a BOOM bug. If $B$ matches one reference
but not the other, the odd reference is the suspect. If the two references
disagree with each other, that is a reference-model problem, apart from BOOM. We
call this triangulation. It turns the hard, human step into a simple vote for
most cases, and only sends a true three-way split to the model. If a bug makes
BOOM run forever, a cycle cap stops it so we still get a full, comparable trace
instead of a timeout with no output.

\section{Implementation}
The loop is about 3.8k lines of Python with no extra runtime libraries. It
drives a toolchain we script fully so it rebuilds the same way on a fresh cloud
account: Spike (pinned to one commit, built from source),
\code{riscv64-unknown-elf-gcc}~14.2, Verilator~5.022, Dromajo, and a Chipyard
BOOM~v3 Verilator model with commit logging on. The whole toolchain installs
without root. The model calls use Gemini and fall back to a fixed offline path
when no key is set, so the loop still runs and can be tested at no cost. Every
result records where it came from: a fixed rule, the offline path, or which
model. The code ships with 58 unit tests, most of them checks that the test
gate and the decision rules reject what they should.

\section{Results}
All numbers below come from runs on the funded compute: BOOM built and run as a
Chipyard \code{MediumBoomV3} Verilator model, with Spike and Dromajo as the two
reference models.

\subsection{Planted-bug study}
Real bugs in a mature core are rare, so we test the loop against planted bugs.
Each planted bug is one small, reversible change to BOOM's source at a chosen
spot, anchored by an exact text match so a wrong match is caught before a build,
not after. The set has eleven planted bugs: three in load/store memory ordering
and eight in forwarding and the datapath. Each one is a known bug, so the test
gives two numbers: \emph{recall}, how many the loop finds, and \emph{precision},
how many of its BOOM-bug calls land on a real planted bug. A third number comes
for free: any BOOM-bug call on the clean version is a false alarm and must be
zero.

Two things keep the numbers honest. We only count recall over bugs that are
\emph{observable}: a bug that never changes committed state cannot be found by
any comparison, so we check each missed bug against the clean baseline and mark
it ``not observable'' instead of counting it as a miss. And we check that each
planted binary really differs from the baseline, so a ``not observable'' result
can never come from a cached or failed build.

\begin{table}[t]
\centering\scriptsize
\setlength{\tabcolsep}{4pt}
\begin{tabular}{lll}
\toprule
Planted bug & Site & Result \\
\midrule
\cls{bypass\_rs1\_disabled}     & register-read & \textbf{found} \\
\cls{bypass\_rs2\_disabled}     & register-read & \textbf{found} \\
\cls{bypass\_rs1\_match\_wrong} & register-read & \textbf{found} \\
\cls{bypass\_rs2\_match\_wrong} & register-read & \textbf{found} \\
\cls{bypass\_rs1\_zero\_guard}  & register-read & \textbf{found} \\
\cls{bypass\_rs2\_zero\_guard}  & register-read & \textbf{found} \\
\cls{forward\_std\_val\_stuck}  & lsu           & \textbf{found} \\
\cls{fwd\_iresp\_early}         & lsu           & \textbf{found} \\
\cls{order\_fail\_suppressed}   & lsu           & not observable \\
\cls{fencei\_early}             & lsu           & not observable \\
\cls{mov\_wrong\_operand}       & exec unit     & not observable \\
\bottomrule
\end{tabular}
\caption{Planted-bug results. Every bug that changes committed state on a single
core was found; the three that never change committed state are left out of
recall, not counted as misses.}
\label{tab:seedresults}
\end{table}

\textbf{Results.} Of the eleven planted bugs (Table~\ref{tab:seedresults}),
eight change committed state under our tests and three do not. The vote found
all eight that do: recall over observable bugs is \textbf{8 of 8}. Precision is
\textbf{100\%}: on the clean baseline the tool made \emph{zero} BOOM-bug calls,
so every call landed on a real planted bug. The three not-observable bugs are
two multi-core memory-ordering effects and one register-move case our tests
never trigger. These are honest limits of single-core bare-metal testing, not
tool failures. Each one builds a binary that really differs from the baseline
yet commits the same state.

\begin{figure}[t]
\centering\scriptsize
\begin{tabular}{@{}ll@{}}
\toprule
\textbf{Spike (reference)} & \textbf{BOOM} \\
\midrule
\code{0x10 mv a0,s2} & \code{0x10 mv a0,s2} \\
\textit{(handler entry)} & \code{0x12 ecall} $\leftarrow$ \textit{logged} \\
\code{0x50 csrr mcause} & \code{0x50 csrr mcause} \\
\bottomrule
\end{tabular}
\caption{Finding 1. BOOM logs the trapping \code{ecall}; Spike does not. Both
reach the same handler with the same \code{mcause}. The apparent control-flow
difference is a logging habit, not a bug.}
\label{fig:finding1}
\end{figure}

\textbf{Finding 1: a logging habit that looks like a bug.} Our exception test
traps on purpose. The comparison first reported a control-flow difference. We
looked instead of assuming. BOOM logs the trapping instruction (for example
\code{ecall}) with no register write, while Spike jumps straight to the handler.
Both then enter the same handler with the same \code{mcause}, so the state is
the same. It is a logging habit, not a bug. Left alone it would raise a false
alarm on every exception test. We handle it, but we \emph{count and report} each
skip instead of quietly realigning, since a tool that quietly realigns can hide
a real bug. Handling it showed a second case, the compressed \code{c.ebreak}
form, that our first list missed. The same habit shows up \emph{between the two
references}: on the exception test Spike and Dromajo disagree, which the vote
flags as a reference-model mismatch. A closer look shows it is the same logging
habit, not a reference-model bug, and the tool does not report it as one.

\textbf{Finding 2: the second reference finds a bug the first misses.} The
\cls{forward\_std\_val\_stuck} bug breaks store-to-load forwarding. The wrong
value throws the program off course, so the difference shows up as a cut-off
trace, which single-reference triage calls a test-setup problem and misses. The
vote finds it: both references finish and agree, BOOM is the odd one out, so the
call is a BOOM bug no matter how the difference looks. This is the clear case for
the second reference: a bug the one-reference path missed, found by a simple
vote.

\textbf{Finding 3: repeat and evidence checks matter.} Two bugs in our own
harness, including a wrong file path that gave an empty BOOM trace, were caught
by the decision part, which correctly refused to call an empty trace a bug. This
is why we run fixed checks before any model call: repeat first, and tell ``the
tool did not run'' apart from ``the models disagree.''

\section{Limitations and Future Work}
\label{sec:limitations}
We are explicit about four limits.

First, single-core bare-metal tests cannot see every kind of bug: three of our
eleven planted bugs never change committed state this way, so no comparison can
catch them. Two are multi-core memory effects that need more than one hart.

Second, our results do not isolate the value of aiming tests at specific RTL
spots. Seven of the eight caught bugs were detected by one broad, mixed program
rather than by the targeted programs, and because the study stops at the first
detection, the targeted programs often did not run. So the loop finds these bugs
by comparing and voting on a datapath-exercising program; whether reading the RTL
to aim tests actually improves detection is not shown here and needs a controlled
comparison. The eight observable bugs are also close relatives (bypass and
forwarding), so the recall number covers a narrow part of the microarchitecture.

Third, the two reference models agreed on all real committed state in this study;
they differed only on the exception-test logging habit, which is not a real
disagreement. So the distinctive value of a second reference---breaking a tie, or
catching a bug in a reference model---is designed in but not yet demonstrated on a
real case. Here the vote gives confidence (both agree and BOOM is the outlier)
rather than resolving a genuine three-way split.

Fourth, these are planted bugs that show the method works; we have not yet found
a new, real bug in BOOM.

\section{Ideas to Extend the Loop}
Three more steps would make the loop stronger. Each builds on existing work.

\textbf{A formal reference model.} Adding the Sail formal RISC-V
model~\cite{sail}, the official executable spec, as a third reference would make
the vote even stronger, since it is grounded in the written rules.

\textbf{Coverage feedback for choosing spots.} Right now the loop reads only
static source. Feeding it live signals, such as CSR changes like
ProcessorFuzz~\cite{processorfuzz} or register coverage like
DIFUZZRTL~\cite{difuzzrtl}, would tell the model where the tests have not been
yet. No current fuzzer uses a model to choose where to aim.

\textbf{Deeper test programs.} Straight-line programs cannot reach every state.
Cascade's~\cite{cascade} tangled programs reach much deeper, which would raise
how many planted bugs are observable and lift the ceiling on recall.

\textbf{Formal check of a difference.} A difference found by testing is a cheap
lead; a bounded formal check is a proof. A found difference could seed a formal
property, checked by the yosys/SymbiYosys flow we built earlier, to turn a
suspect trace into a proof or a small counterexample.

\section{Conclusion}
We built a full, repeatable loop that looks for bugs in BOOM by comparing it to
two reference models. We credit the co-simulation as a Berkeley feature and add
the layer on top: reading the RTL to choose tests, writing tests that do real
work, and deciding what each difference means. The tool is careful not to invent
findings: a gate that drops empty tests, fixed rules that settle the easy cases,
and a high bar for calling something a BOOM bug. On planted bugs it found every
one that a single core can show, made no false alarms, and, with a second
reference model, found a bug the one-reference path missed.

\bibliographystyle{IEEEtran}
\bibliography{refs}

\end{document}
